\section{Introducción y Antecedentes}

\subsection{Agentes del Proyecto}

\subsubsection{Promotor y Adjudicatario del Proyecto}
El promotor de la propuesta de adecuación estética y adjudicatario de la concesión del servicio es la entidad mercantil \textbf{Repostería Barroso Iglesias, SL}, con domicilio social en Sevilla, especialista en gestión de servicios de restauración y alimentación.

\subsubsection{Dirección Técnica y Profesionales Intervinientes}
La redacción de la presente memoria técnica y el diseño de la propuesta de adecuación han sido realizados por el arquitecto y proyectista técnico \textbf{José Ramón Redondo Rosario}, colegiado en Sevilla.

\subsection{Antecedentes y Objeto de la Memoria}
El objeto de esta memoria técnica es definir, justificar y describir detalladamente las actuaciones de remodelación estética, distribución y equipamiento propuestas para la adecuación de los espacios de la cafetería del \textbf{Hospital de Rehabilitación y Traumatología}, perteneciente al \textbf{Hospital Universitario Virgen del Rocío}, de acuerdo con el Pliego de Prescripciones Técnicas (PPT) de la licitación.

El alcance del proyecto contempla exclusivamente la adecuación interior de los comedores (público y de personal) y de la barra de servicio, así como el adecentamiento estético y funcional de la fachada exterior de acceso. Dado que todas las actuaciones son de carácter no estructural (acabados, falsos techos, revestimientos, pintura, iluminación, mobiliario y vegetación), no se altera la volumetría, estructura portante ni la configuración arquitectónica general del edificio.

\subsection{Emplazamiento Objeto de Actuación}
Los espacios objeto de la remodelación se encuentran situados en el complejo hospitalario del Hospital Universitario Virgen del Rocío, concretamente en:
\begin{itemize}
    \item \textbf{Dirección:} Calle Torcuato Luca de Tena s/n, 41013, Sevilla.
    \item \textbf{Edificio:} Hospital de Rehabilitación y Traumatología.
    \item \textbf{Ubicación:} Planta Baja del edificio, con fachada y accesos orientados hacia la vía pública y acera peatonal de la Calle Torcuato Luca de Tena.
\end{itemize}

\subsection{Descripción Geométrica y Superficies Útiles}
La cafetería del Hospital de Rehabilitación y Traumatología cuenta con una superficie útil total de $417,37\text{ m}^2$ y una superficie construida total de $455,43\text{ m}^2$. Los espacios y estancias se adaptan a la geometría preexistente y a los requisitos del servicio sin alterar la distribución estructural de mampostería. En el Cuadro~\ref{tab:superficies} se detalla la relación completa de espacios y sus correspondientes superficies útiles.

\begin{table}[htbp]
    \centering
    \caption{Cuadro de superficies útiles y construidas de la cafetería.}
    \label{tab:superficies}
    \begin{tabular}{lr}
        \toprule
        \textbf{Espacio / Estancia} & \textbf{Superficie Útil ($m^2$)} \\
        \midrule
        Comedor Público & 140,83 \\
        Comedor Personal & 58,28 \\
        Servicio de Barra & 50,03 \\
        Aseo Adaptado & 5,48 \\
        Vestíbulo Aseos & 2,70 \\
        Aseo Masculino & 3,61 \\
        Aseo Femenino & 3,61 \\
        Vestíbulo & 22,47 \\
        Espacio sin Uso & 9,00 \\
        Propano & 2,85 \\
        Climatización & 12,85 \\
        Sala de Lavado & 21,54 \\
        Cocina & 27,10 \\
        Preparaciones & 14,65 \\
        Pasillo Cámaras & 15,00 \\
        Almacén & 7,12 \\
        Cámaras & 14,52 \\
        Cuarto Basuras & 2,49 \\
        Otros (N/A) & 3,24 \\
        \midrule
        \textbf{Superficie Útil Total} & \textbf{417,37} \\
        \textbf{Superficie Construida} & \textbf{455,43} \\
        \bottomrule
    \end{tabular}
\end{table}
